\section{Two variational inputs}\label{sec:trials}

The finite initializations need a lower bound on the untwisted partition
function at lengths $72$ and $120$. One fixed triple of integer matrices
defines periodic matrix-product vectors at both lengths, with energy per bond below
$-a=-700741/500000$. A block decomposition reduces the calculation to
powers of integer matrices of order at most $162$.

\subsection{Norm and energy contractions}

For $N\ge4$ and three real $D$-by-$D$ matrices $A_s$, indexed by $s=-1,0,1$, define
a vector in the standard orthonormal spin-one weight basis by
\begin{equation}\label{eq:trial-vector}
 \psi_N(s_0,\ldots,s_{N-1})
 =\Tr(A_{s_0}\cdots A_{s_{N-1}}).
\end{equation}
This is the periodic matrix-product representation
of~\cite[Section~2.1, Equations~(2)--(3)]{PerezGarcia2007}.
Write
\begin{align}
 B_{st}&=A_sA_t,\qquad T=\sum_{s=-1}^1 A_s\otimes A_s,
 \label{eq:trial-transfer}\\
 O&=\sum_{s,t=-1}^1 st\,B_{st}\otimes B_{st}
   +\sum_{\substack{s=-1,0\\t=0,1}}
    \bigl(B_{st}\otimes B_{s+1,t-1}
       +B_{s+1,t-1}\otimes B_{st}\bigr).
 \label{eq:trial-insertion}
\end{align}
All traces here are ordinary, unnormalized matrix traces.

\begin{lemma}\label{lem:trial-contraction}
Let $N\ge4$ and $D\ge1$ be integers, and let $A_{-1},A_0,A_1\in M_D(\R)$.
Define $\psi_N,T,O$ by Equations~\eqref{eq:trial-vector}--\eqref{eq:trial-insertion}.
For the periodic spin-one Heisenberg Hamiltonian $H_N$, set
\[
 V_N=\Tr(T^N),\qquad U_N=\Tr(T^{N-2}O).
\]
Then
\[
 \norm{\psi_N}^2=V_N,\qquad
 \langle\psi_N,H_N\psi_N\rangle=NU_N.
\]
In particular, if $V_N>0$ and $U_N<-aV_N$, then
$E_0(N)<-Na$ and $Z_N(b)>1$ for every $b>0$.
\end{lemma}

\begin{proof}
Expand $T^N$ and use $\Tr(X\otimes Y)=\Tr(X)\Tr(Y)$.
Since all amplitudes in~\eqref{eq:trial-vector} are real, the result is
the sum of their squares. This proves the norm formula without assuming
that the finite contraction matrix $T$ is self-adjoint or positive.

In the weight basis, a bond has diagonal entry $st$ on the pair $(s,t)$.
Its off-diagonal entries are $1$ between $(s,t)$ and $(s+1,t-1)$
for $s=-1,0$ and $t=0,1$: each follows from
$h=S^z\otimes S^z+(S^+\otimes S^-+S^-\otimes S^+)/2$
and the ladder entries $\sqrt2$.
Inserting this bond into the two amplitude layers gives exactly
$O$ in~\eqref{eq:trial-insertion}. Summing the other physical indices
therefore gives $U_N$ for one bond.
Cyclicity of the amplitude trace makes $\psi_N$ invariant under cyclic
site translation, so every bond, including the periodic seam, contributes
equally.

For $V_N>0$, the variational principle gives
$E_0(N)\le NU_N/V_N<-Na$. The shifted ground energy is negative, so
its term in $\Tr e^{-b(H_N+NaI)}$ exceeds $1$ for every $b>0$.
\end{proof}

\subsection{One integer matrix triple}

Set $D=26$. Index eight types by $\alpha=1,\ldots,8$,
put
\[
 (\ell_1,\ldots,\ell_8)=(1,1,1,1,3,3,3,5),
\]
and order the virtual labels $(\alpha,d)$ first by $\alpha$ and then by
$d=-\ell_\alpha,-\ell_\alpha+2,\ldots,\ell_\alpha$.
The matrices $A_s$ use the diagonal parameters $D_\alpha$ in
Table~\ref{tab:trial-diagonal} and the parameters
$P_{\alpha\beta}$ in Table~\ref{tab:trial-pairs}, defined only when
$\ell_\beta=\ell_\alpha+2$.

\begin{table}[ht]
\centering
\begin{tabular}{c|rrrrrrrr}
$\alpha$&1&2&3&4&5&6&7&8\\ \hline
$D_\alpha$&$-289$&$-84$&$-9$&33&$-22$&27&71&$-11$
\end{tabular}
\caption{Diagonal parameters, common to both lengths.}
\label{tab:trial-diagonal}
\end{table}

\begin{table}[ht]
\centering
\begin{tabular}{c|r@{\qquad}c|r@{\qquad}c|r}
$(\alpha,\beta)$&$P_{\alpha\beta}$&
$(\alpha,\beta)$&$P_{\alpha\beta}$&
$(\alpha,\beta)$&$P_{\alpha\beta}$\\ \hline
$(1,5)$&$-11$&$(2,7)$&$-49$&$(4,6)$&9\\
$(1,6)$&$-35$&$(3,5)$&$-6$&$(4,7)$&$-32$\\
$(1,7)$&$-28$&$(3,6)$&$-25$&$(5,8)$&1\\
$(2,5)$&$-23$&$(3,7)$&$-20$&$(6,8)$&2\\
$(2,6)$&16&$(4,5)$&$-34$&$(7,8)$&$-16$
\end{tabular}
\caption{Parameters joining virtual types whose labels $\ell$ differ by $2$.}
\label{tab:trial-pairs}
\end{table}

Here are complete entry rules. Consider the row $(\alpha,d)$ and column
$(\beta,e)$, put $\ell=\ell_\alpha$, $k=\ell_\beta$, and fix
$s\in\{-1,0,1\}$. The entry is zero unless
\begin{equation}\label{eq:trial-support}
 e=d+2s,\qquad |k-\ell|\le2,\qquad
 (k=\ell\Longrightarrow\alpha=\beta).
\end{equation}
For an allowed entry define
\[
 p=\begin{cases}
 D_\alpha,& k=\ell,\\
 P_{\alpha\beta},& k=\ell+2,\\
 -P_{\beta\alpha},& k=\ell-2.
 \end{cases}
\]
Put $r=(d+\ell)/2$ and set the entry equal to $p f$,
where
\begin{equation}\label{eq:trial-entry}
\begin{array}{c|ccc}
 &s=-1&s=0&s=1\\ \hline
k=\ell+2&
 2(\ell+1-r)(\ell+2-r)&2(r+1)(\ell-r+1)&(r+1)(r+2)\\
k=\ell&
 2(\ell-r+1)&d&-(r+1)\\
k=\ell-2&2&-2&1
\end{array}
\end{equation}
lists the values of $f$.
These entry rules are independent of $N$.

\subsection{Finite contractions and variational bounds}

We now evaluate the contractions in Lemma~\ref{lem:trial-contraction}
for this fixed triple at $N=72$ and $N=120$.
The support rule~\eqref{eq:trial-support}
makes this a short block calculation. On paired virtual labels
$((\alpha,d),(\beta,e))$, both $T$ and $O$ preserve $d-e$.
For $T$, the two layers shift by the same $2s$.
For $O$, the two layers have the same total physical weight $s+t$,
also in each off-diagonal term.
Thus the difference blocks are exhaustive and have dimensions
\begin{equation}\label{eq:trial-blocks}
\begin{array}{c|rrrrrrrrrrr}
d-e&-10&-8&-6&-4&-2&0&2&4&6&8&10\\ \hline
\text{dimension}&1&8&32&80&136&162&136&80&32&8&1
\end{array}.
\end{equation}
They sum to $676=26^2$; each block contributes once.

For completeness, the entire integer computation consists of the
following operations. Form $A_s$, $B_{st}$, $T$, and $O$ from
Tables~\ref{tab:trial-diagonal}--\ref{tab:trial-pairs} and
Equations~\eqref{eq:trial-transfer}, \eqref{eq:trial-insertion},
\eqref{eq:trial-support}, and~\eqref{eq:trial-entry}.
Restrict $T,O$ to each difference block, obtaining $T_\delta,O_\delta$.
Compute powers by
\[
 P_0=I,\qquad
 P_{2j}=P_j^2\quad(j\ge1),\qquad
 P_{2j+1}=P_j^2T_\delta\quad(j\ge0).
\]
Then accumulate
\begin{equation}\label{eq:trial-integer-algorithm}
 U_N=\sum_\delta\sum_{i,j}
       (P_{N-2})_{ij}(O_\delta)_{ji},
 \qquad
 V_N=\sum_\delta\sum_{i,j}
       (P_{N-2})_{ij}(T_\delta^2)_{ji}.
\end{equation}
Only arbitrary-precision integer operations are used.
For $N=72$, the required power is $70$.
For $N=120$, one may use
$T^{118}=T^{64}T^{32}T^{16}T^4T^2$ and
$T^{120}=T^{64}T^{32}T^{16}T^8$ within each block.

The resulting exact integer comparisons are
\begin{equation}\label{eq:trial-integer-results}
 V_N>0,\qquad
 -c_NV_N\le10^{12}U_N<(-c_N+1)V_N,
 \qquad
 \begin{array}{c|r}
 N&c_N\\ \hline
 72&1401482112601\\
 120&1401482105174
 \end{array}.
\end{equation}
The supplied computation files contain the complete matrices and the
full integers $U_N,V_N$, rather than only these shortened intervals.
In particular, the files
\path{verification/computations/evidence/trials/certificate72.json} and
\path{verification/computations/evidence/trials/certificate120.json}
record the full contractions and positive margins
$-700741V_N-500000U_N$. The common matrices are recorded in
\path{verification/computations/evidence/trials/matrix-data.json}.
The program \path{verification/computations/trials/run_trial.py},
with \texttt{-{}-length 72} or \texttt{-{}-length 120}, executes the
entry rules and block contractions in
\path{verification/computations/trials/original_trial.py} at the selected length.
The integer comparisons in~\eqref{eq:trial-integer-results} also
follow directly from their full recorded integers.

\begin{lemma}\label{lem:trial-inputs}
For the untwisted periodic spin-one Heisenberg Hamiltonian and
$a=700741/500000$,
\[
 E_0(72)<-72a,\qquad E_0(120)<-120a.
\]
Consequently $Z_{72}(b)>1$ and $Z_{120}(b)>1$ for every $b>0$.
\end{lemma}

\begin{proof}
For each row of~\eqref{eq:trial-integer-results},
$(-c_N+1)/10^{12}<-700741/500000$.
Thus the corresponding explicit matrices have $V_N>0$ and
$U_N/V_N<-a$. Apply Lemma~\ref{lem:trial-contraction}.
\end{proof}

The support rule puts these trial vectors in total weight zero, but the
variational bound is for the minimum of the full physical spectrum.
Their only symmetry used in the contraction proof is cyclic site
translation. The trial inputs hold at every positive inverse
temperature, so each finite initialization can use them at its own
temperature.
